\begin{abstract}
Filtered vector search combines similarity search over high-dimensional embeddings with structured predicates over metadata. 
The performance of such queries depends not only on predicate selectivity, 
but also on how the structured filter relates to the local geometry of the embedding space. 
When qualifying points are concentrated near the query, independent of it, 
or mostly absent from its neighborhood, 
different execution strategies can lead to different trade-offs between recall and efficiency.

This thesis studies attribute-embedding correlation as a signal for query planning in filtered vector search. 
It introduces a lightweight metric that estimates whether a structured filter is aligned with the local neighborhood of a query in the embedding space. 
The metric is calibrated without manual labels and is used to classify query-filter pairs into correlation regimes that can guide the choice of execution strategy.
\end{abstract}
